\begin{table}[htbp]
\centering
\caption{Respuestas Agregadas a un Shock de 10\% al Precio del Petróleo}
\label{tab:oil_agg}
\begin{threeparttable}
\begin{tabular}{@{}l c c c@{}}
\toprule
Variable & Impacto & Mín/Máx & Trimestre \\
\midrule
PIB & -0.430 & -0.430 & 1 \\
Inflación IPC & +0.190 & +0.190 & 1 \\
Consumo & -0.278 & -0.278 & 1 \\
Empleo & -0.044 & +0.069 & 31 \\
Tipo de Cambio Real & +0.092 & +0.306 & 28 \\
Balanza Comercial & -0.159 & -0.159 & 1 \\
Tasa de Política & +0.122 & +0.128 & 2 \\
\bottomrule
\end{tabular}
\begin{tablenotes}[flushleft]
\footnotesize
\item \textit{Notas:} PIB, consumo, empleo y TCR en \% de desviación del estado estacionario; inflación y tasa de política en pp anualizados; balanza comercial en pp del PIB. El shock es un aumento único de 10\% en el precio mundial del petróleo $P^{O*}_t$ con persistencia $\rho = 0.9$. Impacto = respuesta en el trimestre 1. Mín/Máx = respuesta extrema en 40 trimestres.
\end{tablenotes}
\end{threeparttable}
\end{table}
